%%%%%%%%%%%%%%%%%%%%%%%%%%%%%%%%%%%%%%%%%%%%%%%%%%%%%%%%%%%%%%%%%%%%%%%%%%%%%%%
% Which Divergences Control Representations?
% Tail Sensitivity and the Limits of Identifiability
%
% Target: ICML 2027 (projected abstract Jan 16 / paper Jan 22, 2027 -- CONFIRM
% on icml.cc before relying on these dates).
% Internal freeze: December 31, 2026.
%
% HOW TO USE THIS FILE
%   Each \TODO{} is keyed to a schedule gate (G1 = Oct 5, G2 = Nov 2,
%   G3 = Nov 30). Sections are ordered so that the theory written in weeks 1-8
%   lands directly into Secs. 3-5 as it is proved, rather than being
%   reconstructed in December.
%
%   Proof obligations referenced as [O1]..[O6]:
%     O1 Rate lemma for D_alpha along the rho-family; critical order alpha*
%     O2 Master impossibility: any TV-continuous divergence fails
%     O3 Quantitative lower bound on d_{f,g}, uniform in rho
%     O4 Max-divergence controls representations
%     O5 Conditional control for KL on Theta_delta
%     O6 Tightened constant via Hoffman-Wielandt
%%%%%%%%%%%%%%%%%%%%%%%%%%%%%%%%%%%%%%%%%%%%%%%%%%%%%%%%%%%%%%%%%%%%%%%%%%%%%%%

\documentclass{article}

% --- ICML 2027 style -----------------------------------------------------
% Replace with the official icml2027.sty when released. Until then icml2026.sty
% is layout-compatible; do NOT submit against the wrong year's style file.
\usepackage{microtype}
\usepackage{graphicx}
\usepackage{subfigure}
\usepackage{booktabs}
\usepackage{hyperref}
\usepackage{amsmath,amssymb,amsthm,mathtools}
\usepackage{bm}
\usepackage{xcolor}

% \usepackage{icml2027}            % camera-ready
\usepackage[accepted]{icml2027}    % <-- comment out for anonymous submission

% --- macros --------------------------------------------------------------
\newcommand{\TODO}[1]{\textcolor{red}{\textbf{[TODO: #1]}}}
\newcommand{\GATE}[1]{\textcolor{blue}{\textsf{\small (#1)}}}

\newcommand{\X}{\mathcal{X}}
\newcommand{\Y}{\mathcal{Y}}
\newcommand{\bff}{\mathbf{f}}
\newcommand{\bfg}{\mathbf{g}}
\newcommand{\bfL}{\mathbf{L}}
\newcommand{\bfN}{\mathbf{N}}
\newcommand{\dKL}{d_{\mathrm{KL}}}
\newcommand{\dLLV}{d^{\lambda}_{\mathrm{LLV}}}
\newcommand{\dSVD}{d_{\mathrm{SVD}}}
\newcommand{\dfg}{d_{\bff,\bfg}}
\newcommand{\Dalpha}{D_{\alpha}}
\newcommand{\Dinf}{D_{\infty}}
\newcommand{\psix}{\psi_{x}}
\newcommand{\psiy}{\psi_{y}}
\newcommand{\TV}{\mathrm{TV}}
\newcommand{\Var}{\mathrm{Var}}
\newcommand{\Cov}{\mathrm{Cov}}
\newcommand{\simL}{\sim_{L}}
\newcommand{\Thetadelta}{\Theta_{\delta}}

\theoremstyle{plain}
\newtheorem{theorem}{Theorem}
\newtheorem{lemma}[theorem]{Lemma}
\newtheorem{proposition}[theorem]{Proposition}
\newtheorem{corollary}[theorem]{Corollary}
\theoremstyle{definition}
\newtheorem{definition}[theorem]{Definition}
\newtheorem{assumption}[theorem]{Assumption}
\theoremstyle{remark}
\newtheorem{remark}[theorem]{Remark}
\newtheorem{conjecture}[theorem]{Conjecture}

\icmltitlerunning{Which Divergences Control Representations?}

\begin{document}

\twocolumn[
\icmltitle{Which Divergences Control Representations?\\
           Tail Sensitivity and the Limits of Identifiability}

\icmlsetsymbol{equal}{*}

\begin{icmlauthorlist}
\icmlauthor{Anonymous}{anon}
\end{icmlauthorlist}

\icmlaffiliation{anon}{Anonymous Institution}
\icmlcorrespondingauthor{Anonymous}{anon@example.org}

\icmlkeywords{representational similarity, identifiability, divergences,
              Renyi divergence, representation learning}

\vskip 0.3in
]

\printAffiliationsAndNotice{}

%%%%%%%%%%%%%%%%%%%%%%%%%%%%%%%%%%%%%%%%%%%%%%%%%%%%%%%%%%%%%%%%%%%%%%%%%%%%%%%
\begin{abstract}
%
% DRAFTING NOTE: written against the expected results. Every numeric claim
% carries a \TODO until the corresponding experiment closes. Keep
% "representations" ahead of "divergence" in the opening sentence: the
% divergence family is the instrument, the representations are the object.
%
Whether two networks that induce similar output distributions must learn
similar internal representations depends, we show, on which notion of
distributional closeness is used. For a model family covering autoregressive
language models, contrastive pre-training and supervised classifiers, we ask
which divergences \emph{control} representational dissimilarity, in the sense
of admitting a modulus of continuity that forces representations together as
the divergence vanishes. Indexing divergences by R\'enyi order $\alpha$, we
show that the answer is governed by a sharp threshold. For any two models whose
predictions concentrate on a common label while their tail log-likelihood
ratios diverge, there is a critical order
$\alpha^{\star}=\min_{y}b_{y}/(b_{y}-a_{y})$, determined by the two models'
margin geometry, below which the divergence vanishes and above which it grows
without bound. We prove $\alpha^{\star}>1$ always. Consequently the entire
segment $\alpha\le1$ fails to control representations at once, which covers
total variation, squared Hellinger, Jensen--Shannon and Kullback--Leibler, and
subsumes the known KL counterexample as the boundary case. Above the threshold
control is recovered, with a modulus that degrades both as $\alpha\downarrow1$
and as the models grow confident: we prove
$\sup_{y}|\log(p/p')|\le D_{\alpha}+\log(1/\delta)/(\alpha-1)$ where
$\delta$ lower-bounds the conditional probabilities, so KL itself controls on
the subclass $\Theta_{\delta}$. We further show that $\alpha^{\star}$ cannot be
inflated for free: geometries that raise it provably shrink the
representational separation they were built to hide. This explains why
confident classifiers and language models, which place negligible mass on most
labels, are exactly the regimes in which likelihood agreement carries the least
representational information.
\TODO{fold in headline empirical numbers once C1--C2 close (G3)}
\end{abstract}

%%%%%%%%%%%%%%%%%%%%%%%%%%%%%%%%%%%%%%%%%%%%%%%%%%%%%%%%%%%%%%%%%%%%%%%%%%%%%%%
\section{Introduction}
\label{sec:intro}

\TODO{Write last, after Secs. 3--5 settle. Target 1.25 columns.}

Paragraph plan:
\begin{enumerate}
  \item Representational similarity across independently trained networks is an
        active question; identifiability theory supplies the invariance that a
        similarity measure ought to respect.
  \item Prior work establishes that \emph{exact} likelihood equality pins
        representations down to an invertible linear map, and that small KL
        divergence does not. Frame that as one data point rather than a
        conclusion: the interesting question is which divergences do.
  \item State the organising question: for which $D$ does
        $D \le \epsilon \Rightarrow \dfg \le h(\epsilon)$ with $h(0^{+})=0$?
  \item Answer in three parts: a mechanism (tail insensitivity), an
        impossibility covering a whole topology's worth of divergences, and a
        positive endpoint plus a conditional recovery for KL.
  \item Practical consequence: model confidence, not model quality, governs how
        much representational information likelihood agreement carries.
\end{enumerate}

\paragraph{Contributions.}
\begin{itemize}
  \item A definition of \emph{control} (Def.~\ref{def:control}) that makes both
        positive and negative claims about divergence--representation coupling
        falsifiable and uniform over a model class.
  \item An impossibility result covering all divergences continuous in total
        variation (Thm.~\ref{thm:master}), subsuming the known KL failure, with
        a quantitative lower bound on the residual dissimilarity
        (Thm.~\ref{thm:lower}).
  \item A critical R\'enyi order $\alpha^{\star}$ (Thm.~\ref{thm:alphastar})
        determined by tail geometry, with matching empirics.
  \item Positive results: max-divergence controls
        (Thm.~\ref{thm:maxdiv-control}); KL controls on $\Thetadelta$
        (Thm.~\ref{thm:kl-conditional}); a tightened constant improving the
        known bound by $\sqrt{M}$ (Prop.~\ref{prop:tighten}).
  \item Experiments on constructed models, synthetic angular classification and
        CIFAR-10 confirming the mechanism, including a label-smoothing
        intervention that \emph{restores} KL's predictive power as predicted.
\end{itemize}

%%%%%%%%%%%%%%%%%%%%%%%%%%%%%%%%%%%%%%%%%%%%%%%%%%%%%%%%%%%%%%%%%%%%%%%%%%%%%%%
\section{Setup}
\label{sec:setup}

\subsection{Model class and linear identifiability}
\label{sec:model-class}

Models are pairs $(\bff,\bfg)\in\Theta$ with $\bff:\X\to\mathbb{R}^{M}$ the
embedding and $\bfg:\Y\to\mathbb{R}^{M}$ the unembedding, inducing
\begin{equation}
  p_{\bff,\bfg}(y\mid x)
  = \frac{\exp(\bff(x)^{\top}\bfg(y))}
         {\sum_{y'\in\Y}\exp(\bff(x)^{\top}\bfg(y'))}.
  \label{eq:model}
\end{equation}
Displaced maps $\bff_{0}(x)=\bff(x)-\bff(x_{0})$,
$\bfg_{0}(y)=\bfg(y)-\bfg(y_{0})$; $\bfL$ has columns $\bfg_{0}(y)$ and $\bfN$
columns $\bff_{0}(x)$.

\TODO{Compress prior setup to $\approx$0.6 column. Restate the diversity
condition and linear identifiability as background, cite, do not reprove.}

\subsection{Representational dissimilarity}
\label{sec:repdis}

\TODO{Recall $\dSVD$ and $\dfg$. State the invariance class in one sentence and
defer the comparison with CCA-based measures to Sec.~\ref{sec:discussion}.}

%%%%%%%%%%%%%%%%%%%%%%%%%%%%%%%%%%%%%%%%%%%%%%%%%%%%%%%%%%%%%%%%%%%%%%%%%%%%%%%
\section{When Does a Divergence Control Representations?}
\label{sec:control}

\begin{definition}[Control]
\label{def:control}
A divergence $D$ on conditional distributions \emph{controls} $\dfg$ on a model
class $\Theta$ if there exists a nondecreasing $h:\mathbb{R}_{+}\to\mathbb{R}_{+}$
with $h(0^{+})=0$ such that for all
$(\bff,\bfg),(\bff',\bfg')\in\Theta$,
\begin{equation}
  D\!\left(p_{\bff,\bfg},p_{\bff',\bfg'}\right)\le\epsilon
  \;\Longrightarrow\;
  \dfg\!\left((\bff,\bfg),(\bff',\bfg')\right)\le h(\epsilon).
  \label{eq:control}
\end{equation}
We call $h$ a \emph{modulus of control} for $D$ on $\Theta$.
\end{definition}

\begin{remark}
The quantifiers matter. Control is uniform over $\Theta$; a divergence may fail
on $\Theta$ yet control on a subclass, which is exactly the situation for KL
(Sec.~\ref{sec:positive}).
\end{remark}

\TODO{One short paragraph: recover the known results as instances. Exact
likelihood equality is $h\equiv 0$; the LLV bound is $h(\epsilon)=2M\epsilon$;
the known KL counterexample is a failure of \eqref{eq:control}. \GATE{G1}}

%%%%%%%%%%%%%%%%%%%%%%%%%%%%%%%%%%%%%%%%%%%%%%%%%%%%%%%%%%%%%%%%%%%%%%%%%%%%%%%
\section{Tail Insensitivity and Impossibility}
\label{sec:negative}

\subsection{The co-concentration mechanism}
\label{sec:mechanism}

The key structural fact about the norm-growth construction is that \emph{both}
models converge to the same point mass while their tail log-ratios diverge.

\begin{lemma}[Co-concentration]
\label{lem:coconc}
Along the $\rho$-indexed family, for every $x\in\X$,
$p_{\bff,\bfg}(\hat y\mid x)\to1$ and $p_{\bff',\bfg'}(\hat y\mid x)\to1$ for a
\emph{common} $\hat y$, hence
$\|p_{\bff,\bfg}(\cdot\mid x)-p_{\bff',\bfg'}(\cdot\mid x)\|_{\TV}\to0$, while
$\sup_{y}|\log p_{\bff,\bfg}(y\mid x)-\log p_{\bff',\bfg'}(y\mid x)|=\Theta(\rho)$.
\end{lemma}
\begin{proof}
\TODO{[O2] Short; assemble from the existing concentration argument plus the
angle-matching assumption. \GATE{G1}}
\end{proof}

\begin{theorem}[Master impossibility]
\label{thm:master}
Let $D$ be any divergence that vanishes along sequences converging in total
variation. Then $D$ does not control $\dfg$ on $\Theta$. In particular total
variation, squared Hellinger, Jensen--Shannon, every bounded $f$-divergence,
and Kullback--Leibler all fail to control $\dfg$.
\end{theorem}
\begin{proof}
\TODO{[O2] Immediate from Lem.~\ref{lem:coconc} plus
Thm.~\ref{thm:lower}. KL needs its own line since it is not TV-continuous in
general; cite the existing rate. \GATE{G1}}
\end{proof}

\subsection{How far apart do the representations stay?}
\label{sec:lower}

\begin{theorem}[Uniform separation]
\label{thm:lower}
There is $c>0$, independent of $\rho$ and explicit in the cluster geometry,
such that $\max\dSVD\ge c$ for every member of the family.
\end{theorem}
\begin{proof}
\TODO{[O3] The rotation is explicit ($\pi/2$ in the first two coordinates), so
compute the cross-covariance singular values in closed form. This upgrades
``not linearly equivalent'' to ``quantitatively far'' and is what makes the
impossibility bite. \GATE{G2}}
\end{proof}

\subsection{A critical R\'enyi order}
\label{sec:alphastar}

Fix an input $x$ and write $s_{y}=\cos(\bff(x),\bfg(y))$, $\hat y=\arg\max_{y}s_{y}$,
and let
\begin{equation}
  a_{y}:=s_{\hat y}-s_{y}>0, \qquad b_{y}:=s'_{\hat y}-s'_{y}>0
  \label{eq:gaps}
\end{equation}
be the margin gaps under the two models. Logits are $c\rho s_{y}$, so
$p(y\mid x)\propto e^{-c\rho a_{y}}$ and $p'(y\mid x)\propto e^{-c\rho b_{y}}$.
Define the \emph{tilted exponent}
\begin{equation}
  E_{\alpha}(y) \;:=\; \alpha a_{y}+(1-\alpha)b_{y} \;=\; b_{y}+\alpha(a_{y}-b_{y}).
  \label{eq:tilted}
\end{equation}

\begin{theorem}[Rate lemma and critical order]
\label{thm:alphastar}
Let $\mathcal{P}:=\{y\neq\hat y: a_{y}\neq b_{y}\}$ and suppose
$\mathcal{P}\neq\emptyset$. Put
\begin{equation}
  \alpha^{\star}\;:=\;\min_{y\,:\,b_{y}>a_{y}}\frac{b_{y}}{\,b_{y}-a_{y}\,}.
  \label{eq:alphastar}
\end{equation}
Then $\alpha^{\star}>1$, and as $\rho\to\infty$:
\begin{enumerate}
  \item[(i)] for $\alpha<\alpha^{\star}$, $\ \Dalpha=\Theta\!\big(e^{-c\rho\kappa(\alpha)}\big)$,
        where $\kappa(\alpha)$ is the smallest value in
        $\{E_{\alpha}(y)\}\cup\{a_{y}\}\cup\{b_{y}\}$ carrying a non-vanishing
        total coefficient (labels with $a_{y}=b_{y}$ cancel exactly and are
        excluded);
  \item[(ii)] at $\alpha=1$ all first-order coefficients cancel and
        $\dKL=\Theta\!\big(\rho\,e^{-c\rho\min_{\mathcal{P}}a_{y}}\big)$;
  \item[(iii)] for $\alpha>\alpha^{\star}$, $\ \Dalpha$ grows \emph{linearly} in
        $\rho$ with slope $\max_{y}(-E_{\alpha}(y))/(\alpha-1)$, increasing to
        $\Dinf/\rho\to\max_{y}(b_{y}-a_{y})$ as $\alpha\to\infty$.
\end{enumerate}
\end{theorem}
\begin{proof}
\TODO{[O1] Write up. Expand $\log\sum_{y}p^{\alpha}p'^{1-\alpha}$ to first order
and group the three exponent families by value; $\alpha^{\star}>1$ is immediate
from $0<b_{y}-a_{y}<b_{y}$. \GATE{G1 -- VERIFIED NUMERICALLY 2026-09-22:
crossover matched to $1.5\times10^{-5}$ relative, exponential rates to $<2.3\%$,
linear slopes to $<0.1\%$; see \texttt{trackA/perx.py}.}}
\end{proof}

\begin{corollary}[The classical divergences all sit below threshold]
\label{cor:below}
Since $\alpha^{\star}>1$, the family defeats every divergence of order
$\alpha\le1$. In particular total variation, squared Hellinger,
Jensen--Shannon and Kullback--Leibler fail to control $\dfg$, recovering the
known KL result as the boundary case $\alpha=1$ of a one-parameter statement.
\end{corollary}

\begin{proposition}[The threshold cannot be inflated for free]
\label{prop:tradeoff}
$\alpha^{\star}$ is large only when $a_{y}\approx b_{y}$ for the critical
label, which simultaneously drives the representational signal
$\max_{y}(b_{y}-a_{y})$ toward zero. Hence constructions that defeat
higher-order divergences are exactly those whose representations are less
dissimilar.
\end{proposition}
\begin{proof}
\TODO{[O1b] Make the trade-off quantitative: lower-bound
$\max_{y}(b_{y}-a_{y})$ by a decreasing function of $\alpha^{\star}$.
Empirically $\mathrm{corr}=-0.59$ over 360 geometries, with mean signal
$0.89$ in the top $\alpha^{\star}$ quintile against $1.60$ in the bottom.
This proposition is what upgrades the negative result into an explanation of
the positive one; prioritise it. \GATE{G2}}
\end{proof}

\begin{remark}[Monotonicity gives a well-posed threshold]
\label{rem:monotone}
$\Dalpha$ is nondecreasing in $\alpha$, so by the transfer lemma the set of
controlling orders is upward closed and
$\alpha^{\star}_{\Theta}:=\inf\{\alpha:\Dalpha\text{ controls }\dfg\text{ on }\Theta\}$
is well defined. Theorem~\ref{thm:alphastar} lower-bounds it;
Thm.~\ref{thm:maxdiv-control} with \eqref{eq:suff} upper-bounds it on
$\Thetadelta$.
\end{remark}

\begin{figure}[t]
\centering
% \includegraphics[width=\columnwidth]{fig/ladder.pdf}
\caption{\textbf{Divergence ladder along the norm-growth family.} Each curve is
one divergence as the unembedding norm $\rho$ grows; representational
dissimilarity (dashed) is flat throughout. Orders below $\alpha^{\star}$
collapse toward zero; the max-divergence and $\dLLV$ do not.
\TODO{from A1 \GATE{G1}}}
\label{fig:ladder}
\end{figure}

\begin{figure}[t]
\centering
% \includegraphics[width=\columnwidth]{fig/alphastar.pdf}
\caption{\textbf{Predicted vs.\ measured critical order.} Each point is one
permutation geometry; $x$-axis is $b_{y}/(b_{y}-a_{y})$ from
\eqref{eq:alphastar}, $y$-axis the empirical crossover. Proximity to the
diagonal tests the quantitative content of Thm.~\ref{thm:alphastar}.
\TODO{from A3 \GATE{G2}}}
\label{fig:alphastar}
\end{figure}

%%%%%%%%%%%%%%%%%%%%%%%%%%%%%%%%%%%%%%%%%%%%%%%%%%%%%%%%%%%%%%%%%%%%%%%%%%%%%%%
\section{Which Divergences Do Control}
\label{sec:positive}

\subsection{Max-divergence}
\label{sec:maxdiv}

\begin{theorem}[Max-divergence controls]
\label{thm:maxdiv-control}
Under the invertibility and non-degeneracy conditions of
Sec.~\ref{sec:setup}, if
$\Dinf = \sup_{x,y}\left|\log p_{\bff,\bfg}(y\mid x)-\log p_{\bff',\bfg'}(y\mid x)\right|\le\epsilon$
then
\begin{equation}
  \dfg \;\le\; C(M)\,\frac{\epsilon}{\psi_{\min}},
  \qquad \psi_{\min}:=\min_{i}\psix(y_{i};p).
  \label{eq:maxdiv-bound}
\end{equation}
\end{theorem}
\begin{proof}
\TODO{[O4] Route through the Lemma-4.1 error term; a uniform sup bound on log
ratios bounds each entry directly. \GATE{G2}}
\end{proof}

\begin{lemma}[R\'enyi dominates the sup log-ratio up to a confidence penalty]
\label{lem:suff}
For $\alpha>1$, let $r_{y}:=\log p(y\mid x)-\log p'(y\mid x)$ and
$\delta:=\min_{x,y}p(y\mid x)$. Then
\begin{equation}
  \max_{y} r_{y} \;\le\; \Dalpha(p\|p') \;+\; \frac{\log(1/\delta)}{\alpha-1}.
  \label{eq:suff}
\end{equation}
\end{lemma}
\begin{proof}
Bound the sum by its largest term and use
$\alpha\log p+(1-\alpha)\log p'=\log p+(\alpha-1)r_{y}$:
$\Dalpha\ge\frac{1}{\alpha-1}\max_{y}[\log p(y)+(\alpha-1)r_{y}]
\ge\max_{y}r_{y}-\frac{\log(1/\delta)}{\alpha-1}$.
\TODO{[O4] Symmetrise over the two orderings and chain into
Thm.~\ref{thm:maxdiv-control}. Verified numerically; the bound is tight to
$\approx6\%$ at $\alpha=20$. \GATE{G2}}
\end{proof}

\begin{remark}[The two axes are one]
\label{rem:unify}
Inequality \eqref{eq:suff} is the paper's organising identity: control requires
$\alpha-1\gtrsim\log(1/\delta)$, so the R\'enyi order and the confidence floor
trade off directly. Setting $\alpha\to\infty$ gives
Thm.~\ref{thm:maxdiv-control}; holding $\alpha=1$ and bounding $\delta$ gives
Thm.~\ref{thm:kl-conditional}. The two positive results are endpoints of one
statement, not separate phenomena.
\end{remark}

\begin{remark}[Why LLV normalises by $\psi$]
\label{rem:whypsi}
The $\psi_{\min}$ in \eqref{eq:maxdiv-bound} is model-dependent, so
$\Dinf$ alone yields a non-uniform modulus. Normalising the log-ratios by
$\psix,\psiy$ removes exactly this dependence, which identifies $\dLLV$ as the
scale-corrected member of the same family rather than an unrelated
construction.
\TODO{This remark is load-bearing for the narrative; write it carefully.}
\end{remark}

\subsection{Conditional control for KL}
\label{sec:klcond}

\begin{theorem}[KL controls on $\Thetadelta$]
\label{thm:kl-conditional}
Let $\Thetadelta=\{(\bff,\bfg)\in\Theta:\min_{x,y}p_{\bff,\bfg}(y\mid x)\ge\delta\}$.
Then $\dKL$ controls $\dfg$ on $\Thetadelta$ with modulus
$h_{\delta}(\epsilon)=C(M)\sqrt{\epsilon/\delta}$.
\end{theorem}
\begin{proof}
\TODO{[O5] Mass weighting is bounded below by $\delta$, so KL dominates the
unweighted log-ratio discrepancy up to $\delta$. Verify the exponent on
$\epsilon$. \GATE{G2}}
\end{proof}

\begin{corollary}[Confidence, not quality]
\label{cor:confidence}
Two models may have arbitrarily small $\dKL$ and arbitrarily large $\dfg$
precisely when $\delta$ is small. Since accurate classifiers and language models
place negligible mass on most labels, these are the regimes in which likelihood
agreement carries the least representational information.
\end{corollary}

\subsection{A tighter constant}
\label{sec:tighten}

\begin{proposition}[$\sqrt{M}$ improvement]
\label{prop:tighten}
Replacing the per-singular-value Weyl bound with Hoffman--Wielandt and applying
Cauchy--Schwarz to the mean yields
$\tfrac1M\sum_{i}|\sigma_{i}-\tilde\sigma_{i}|\le\|E\|_{F}/\sqrt{M}$, improving
the modulus from $2M\epsilon$ to $O(\sqrt{M}\epsilon)$ and widening the
non-vacuous range to $\epsilon<O(1/\sqrt{M})$.
\end{proposition}
\begin{proof}
\TODO{[O6] Also revisit the entrywise covariance inequality on
$E_{ij}=\Cov[z_{i},z_{j}-w_{j}]$, which assumes simultaneous perfect
correlation across all $M^{2}$ entries and is likely loose by a further factor.
\GATE{G2}}
\end{proof}

%%%%%%%%%%%%%%%%%%%%%%%%%%%%%%%%%%%%%%%%%%%%%%%%%%%%%%%%%%%%%%%%%%%%%%%%%%%%%%%
\section{Experiments}
\label{sec:experiments}

We report (i) constructed-model tests of the mechanism, (ii) synthetic angular
classification across widths and class counts, and (iii) CIFAR-10, including an
intervention that manipulates $\delta$ directly. Full protocols in
App.~\ref{app:expdetails}; code at \TODO{repo URL}.

\paragraph{Primary metric: control curves.}
For each divergence we plot the empirical modulus
$\hat h(\epsilon)=\sup\{\dfg : D\le\epsilon\}$, the direct counterpart of
Def.~\ref{def:control}. Failure of $\hat h$ to approach zero as
$\epsilon\to0$ is evidence of non-control.

\paragraph{Statistics.}
All intervals are 95\% cluster bootstrap over \emph{seeds}, not over pairs:
$n$ seeds give $\binom{n}{2}$ dependent pairs, so pair-level resampling is
anticonservative. \TODO{state $n$ per setting}

\subsection{Constructed models: the mechanism}
\label{sec:exp-constructed}
\TODO{A1--A3. Fig.~\ref{fig:ladder}, Fig.~\ref{fig:alphastar}, and a rate table
($\log\Dalpha$ vs.\ $\rho$ slopes against theory). \GATE{G1--G2}}

\subsection{Synthetic angular classification}
\label{sec:exp-synthetic}
\TODO{B1 control curves; B2 confidence intervention via temperature and label
smoothing; B3 width sweep \emph{at matched training loss} (the unmatched
comparison confounds inductive bias with fit quality -- report the regression
with achieved loss as covariate); B4 scaling in class count. \GATE{G3}}

\begin{figure*}[t]
\centering
% \includegraphics[width=\textwidth]{fig/control_curves.pdf}
\caption{\textbf{Control curves.} Empirical modulus $\hat h(\epsilon)$ per
divergence, across seed pairs. Panels: synthetic (4/6/10/18 classes) and
CIFAR-10 ($M=2,3,5$). \TODO{from B1, C1 \GATE{G3}}}
\label{fig:control}
\end{figure*}

\subsection{CIFAR-10}
\label{sec:exp-cifar}
\TODO{C1 tail localisation: divergences restricted to top-$k$ vs.\ tail,
$k=1..9$; prediction is that KL saturates by $k=1$ while $\dfg$ is driven by
$k\ge3$. C2 label-smoothing intervention (the causal test; prioritise). C3
non-vacuity fraction per $M$ under Prop.~\ref{prop:tighten}. C4 robustness
across $m_{\mathrm{CCA}}$, linear CKA, Procrustes, $\dfg$ -- keep in main text,
reviewers will ask. \GATE{G3}}

\subsection{Language modelling}
\label{sec:exp-lm}
\TODO{Optional (Track D). Small tied-embedding transformer, divergences on a
subsampled label set. Drop if week 10 is tight; the paper stands without it,
but it closes the gap between the stated motivation and the evidence.}

%%%%%%%%%%%%%%%%%%%%%%%%%%%%%%%%%%%%%%%%%%%%%%%%%%%%%%%%%%%%%%%%%%%%%%%%%%%%%%%
\section{Related Work}
\label{sec:related}
\TODO{Four short paragraphs, $\approx$0.75 column: (i) representational
similarity measures and their invariances; (ii) linear identifiability for this
model family; (iii) robustness of identifiability under approximate likelihood
maximisation, including $\varepsilon$-non-identifiability; (iv) R\'enyi and
max-divergence outside this setting, where the max-divergence is already the
standard tool for worst-case ratio control. Point (iv) is a genuine connection,
not decoration.}

%%%%%%%%%%%%%%%%%%%%%%%%%%%%%%%%%%%%%%%%%%%%%%%%%%%%%%%%%%%%%%%%%%%%%%%%%%%%%%%
\section{Discussion}
\label{sec:discussion}

\paragraph{Limitations.}
\TODO{Be direct. (a) Control is stated for embeddings and unembeddings, not
intermediate layers. (b) $\dfg$ has a specific invariance class, narrower than
CCA-based measures; we do not claim it is the right one. (c) The label-set
diversity assumption limits applicability when classes are fewer than the
representation dimension. (d) $\alpha^{\star}$ is derived for a particular
adversarial family; whether trained models realise comparable geometries is an
empirical question we only partly answer. Do not bury (d).}

\paragraph{What this changes in practice.}
\TODO{One paragraph. Report $\delta$ alongside any KL-based similarity claim;
prefer high-order or max-divergence when comparing confident models; treat
matched loss as necessary but far from sufficient for matched representations.}

%%%%%%%%%%%%%%%%%%%%%%%%%%%%%%%%%%%%%%%%%%%%%%%%%%%%%%%%%%%%%%%%%%%%%%%%%%%%%%%
\section*{Impact Statement}
\TODO{Standard ICML paragraph. This is foundational work on when representation
comparisons are meaningful; the main societal relevance is that overclaiming
representational agreement from likelihood agreement can mislead
interpretability and distillation pipelines.}

\section*{Acknowledgements}
\TODO{Omit for anonymous submission.}

\bibliographystyle{icml2027}
\bibliography{references}

\newpage
\appendix
\onecolumn

\section{Notation}
\label{app:notation}
\TODO{Table, extended from the prior paper with $\Dalpha,\Dinf,\alpha^{\star},
\Thetadelta,\delta,h(\cdot)$.}

\section{Proofs for Section~\ref{sec:negative}}
\label{app:neg}
\TODO{[O1] rate lemma; [O2] master impossibility; [O3] uniform separation.
\textbf{Write these in weeks 1--4 as they are proved, not in December.}}

\section{Proofs for Section~\ref{sec:positive}}
\label{app:pos}
\TODO{[O4] max-divergence; [O5] conditional KL; [O6] tightened constant.
Weeks 5--8.}

\section{Experimental Details}
\label{app:expdetails}
\TODO{Constructed-model protocol; synthetic data generation; CIFAR-10 training
(ResNet18 embedding, $M\in\{2,3,5\}$, 10 seeds, smoothing arms); divergence
estimation and the $\X_{\mathrm{LLV}}/\Y_{\mathrm{LLV}}$ selection procedure --
state plainly that the reported value is a minimum over a random search and
report sensitivity to the number of draws.}

\section{Additional Results}
\label{app:extra}
\TODO{Per-$M$ and per-class-count breakdowns; alternative similarity measures;
compute budget.}

\end{document}
